\chapter{Refinement and Locality}

\section{Local Refinement Operators}

Locality restricts a refinement rule to a bounded neighborhood, bounded dependency set, or bounded observation window.

\begin{definition}[Local refinement]
A refinement rule \(R_t\) is local at radius \(r\) if the emitted transcript \(T_t\) depends only on data contained in the radius-\(r\) neighborhood selected by the rule.
\end{definition}

\section{Normalization Principle}

A normalization principle converts a broad process into a controlled sequence of local refinements.  The normalized form is useful only when the conversion preserves the obstruction being studied.

\begin{lemma}[Locality-preserving normalization]
If a process admits a locality-preserving normalization and each normalized step has capacity at most \(C\), then a \(T\)-step normalized computation exposes at most \(TC\) units of transcript information.
\end{lemma}

\begin{proof}
Combine the normalization hypothesis with the finite transcript budget from the foundations chapter.
\end{proof}

\section{Boundary}

This chapter supplies a reusable interface.  It does not prove that every algorithm, physical process, or geometric evolution admits such a normalization.
